% Bibliography (Chinese edition).  Only verified entries: author(s), year, title, publisher or
% journal, and a DOI or stable link where one exists.  Chapters add their sources as they are
% written; candidate works listed in docs/PLAN.md must be checked before they are added here.
\begin{thebibliography}{99}
\addcontentsline{toc}{chapter}{参考文献}
\small\raggedright\setlength{\itemsep}{5pt}\setlength{\parskip}{0pt}

\bibitem{devroey} M. De Vroey. \emph{A History of Macroeconomics from Keynes to Lucas and Beyond}. Cambridge University Press, 2016. 全书思想史部分的二手参照。

\bibitem{snowdonvane} B. Snowdon and H. R. Vane. \emph{Modern Macroeconomics: Its Origins, Development and Current State}. Edward Elgar, 2005. 各流派观点与访谈的二手参照。

% ---- 史料（第 12 章起）。数字文本：殆知阁古代文献 daizhigev20，https://github.com/garychowcmu/daizhigev20 ，
% 提交 4a6d6f2088825f132521d848c2ea86cf9c9a7620；引文逐字核对见 data/fetch/verify_quotes.py。
\bibitem{songchaoshishi} 李攸.《宋朝事实》卷十五. 数字文本据殆知阁古代文献（\url{https://github.com/garychowcmu/daizhigev20}，提交 4a6d6f2）。

\bibitem{songshi} 脱脱等.《宋史》（1345 年成书），卷一百八十一《食货志下三》及列传. 数字文本据殆知阁古代文献（提交 4a6d6f2）；卷一百八十一另与维基文库本比对。

\bibitem{yuanshi} 宋濂等.《元史》（1370 年成书），卷九十三《食货志一》、卷九十七《食货志五》、卷二十三与卷二十四（本纪）. 数字文本据殆知阁古代文献（提交 4a6d6f2）。

\bibitem{mingshi} 张廷玉等.《明史》（1739 年刊行），卷八十一《食货志五》. 数字文本据殆知阁古代文献（提交 4a6d6f2）。

\bibitem{xuwenxiantongkao} 《钦定续文献通考》卷九. 数字文本据殆知阁古代文献（提交 4a6d6f2）。

% ---- 研究文献（第 12 章起）
\bibitem{vonglahn1996} R. von Glahn. \emph{Fountain of Fortune: Money and Monetary Policy in China, 1000--1700}. University of California Press, 1996.

\bibitem{pengxinwei} 彭信威.《中国货币史》. 上海人民出版社, 1958. 英译本：Peng Xinwei, \emph{A Monetary History of China}, trans. E. H. Kaplan, 2 vols. Western Washington University, 1994.

\bibitem{guan2024} H. Guan, N. Palma and M. Wu. The rise and fall of paper money in Yuan China, 1260--1368. \emph{Economic History Review}, 77(4): 1222--1250, 2024. \url{https://doi.org/10.1111/ehr.13305}.

\bibitem{jackson1990} P. Jackson. ČĀV. \emph{Encyclopaedia Iranica}, vol. V, fasc. 1, pp. 96--97, 1990. \url{https://www.iranicaonline.org/articles/cav}.

\bibitem{cagan1956} P. Cagan. The monetary dynamics of hyperinflation. In M. Friedman (ed.), \emph{Studies in the Quantity Theory of Money}, pp. 25--117. University of Chicago Press, 1956.

\end{thebibliography}
